% =============================================================================
%  FRAGMENT: Upper Volume, Articles 42--66
%  Source: translation/pages_43-46.md through pages_55-58.md
% =============================================================================

% --- Article 42 ---
\article{42}{When the Market Shows No Movement and All Ten out of Ten Grow Weary}

\begin{sokyubox}
{\small\itshape\color{sokyubar} \jp{第四十二章\quad 相塲無高下十人が十人退屈の事}}

\medskip
When the market shows no significant movement for two or three
months, or merely oscillates within a range-bound market
(\jt{通相塲}{kayoi soba}), all ten out of ten grow weary. Even the
bullish turn bearish, while sellers, convinced their strategy is
working, pile on still more short positions. Thereafter the market
invariably rises. At that moment, bulls and bears alike realize the
truth at once and scramble to buy in a single rush, sending prices
leaping upward in a sharp advance. When all ten out of ten lean to
one side, the market will without fail go the opposite way. If things
came as one expected, trading would be easy---but they do not come as
expected, and calculation cannot reach that far. This is the natural
principle of yin and yang (\jt{陰陽}{in-yo}).
\end{sokyubox}

\commentary
This chapter teaches that when all ten out of ten grow weary, the
moment of change is near. Further, when all ten out of ten lean to one
side, bullish or bearish, the market moves contrary to what everyone
expects; the chapter offers illustrative proof of this principle.

When the market has shown no great movement for two or three months, or
has merely traded back and forth within a range, the crowd grows weary.
Those who had been buying on bullish conviction, seeing the market
unexpectedly fail to rally, lose heart and turn bearish. The sellers,
meanwhile, seeing the market look weak, ride their advantage and pile on
yet more short positions. But thereafter the market rises contrary to
all expectations. At that point, people exclaim: ``So that was the floor
after all---we were fooled by sentiment---one cannot let one's guard
down in this market!'' Bulls and bears alike unite and scramble to buy
all at once, and the market breaks away in a sharp advance. In general,
when all ten out of ten lean to one side, the market will invariably
move in the opposite direction; there is no need for deep
deliberation---one should simply take a position on the opposite side
from the crowd. This is what is called the natural principle of yin-yang
reversion (\jt{陰陽轉還}{in-yo tenkan}).

To explain this principle of yin-yang reversion in somewhat more
scholarly terms, it is as follows.

The principle by which a ceiling is formed consists of four elements:
(1)~buying power is exhausted; (2)~profit-taking (\jt{利喰}{rigui})
selling emerges; (3)~the sprouts of bearish factors become visible; and
(4)~new selling appears. In every case, a price high enough to
constitute a ceiling means that the bulls have already bought to the
utmost limit of their capacity, and it is this full commitment that
produced the extreme high. Since there is no longer any surplus buying
power, profit-taking selling must inevitably emerge. As this
profit-taking selling continues and the market begins to stall, bearish
factors that had been hidden now come to the surface, inviting new
selling; moreover, sellers who had been waiting for a high price add
their weight, and a natural downward momentum takes hold. The previous
high becomes the ceiling, and the market turns steadily lower. The above
is an explanation of the principle by which yang reverts to yin. Next,
the principle by which yin reverts to yang:

The principle by which a floor is formed consists of four elements:
(1)~selling power is exhausted; (2)~profit-taking buy-backs
(\textit{rigui umemodoshi}) emerge; (3)~bullish factors begin to
sprout; and (4)~new buying appears. In general, the sellers have been
riding the momentum of months of victories, increasing their positions
as far as their resources allow and hammering the market down; it is
this that produced the floor. Since they can no longer increase their
positions further, they must perforce begin profit-taking buy-backs. And
when such buy-backs cause the market to recover slightly, some cry ``So
the earlier low was the floor after all!'' and rush to take profits,
while others seize upon some bullish factor or another and new buying
emerges. At times, the urgency of profit-taking buy-backs produces the
paradox known as ``rice sellers cause a rice rally''
(\jp{米売りの米上げ}).

What has been explained above is the scholarly account of the principle
of yin-yang reversion. The ancients, however, did not explain this
principle clearly; not a few described it in mystical or wondrous terms.
Mokko-ken and Jiunsai (\jp{慈雲齋}) are of this sort---a case of ``the
hero deceives people,'' as the saying goes. Jiunsai wrote:

\begin{quote}
\textit{``Between reason and unreason lies a reason beyond reason---know
this as the source of the rise and fall of rice.''}

\textit{``The principle of high and low is an empty principle, invisible
to the eye; that which has neither shadow nor form is its true
substance.''}
\end{quote}

And further:

\begin{quote}
\textit{``The Great Ultimate moves and begets yang; movement reaches its
extreme and becomes stillness. Stillness begets yin; stillness reaches
its extreme and again there is movement. One movement, one
stillness---herein lies the root.''}
\end{quote}

Mokko-ken likewise wrote:

\begin{quote}
\textit{``When the spirit of man has risen to the extreme of yang, yin
arrives and it tends downward. When it has fallen to the extreme of yin,
yang arrives and it rises.''}
\end{quote}

Such pronouncements affect an air of profundity, but when explained in
scholarly terms as above, they amount to nothing remarkable. Yet yin-yang
reversion is not merely a truth of the market world alone---it is a
universal truth of the cosmos. We do not in the least belittle it; but
because such explanations are excessively tinged with a religious
flavor, we fear they may mislead those who come after, and so we add
this word of clarification.

% --- Article 43 ---
\article{43}{How to Buy the November Contract}

\begin{sokyubox}
{\small\itshape\color{sokyubar} \jp{第四十三章\quad 霜月限買入心得方の事}}

\medskip
For the November contract (\jt{霜月限}{shimotsuki-giri}), buying as
early as the sixth or seventh month leads only to tedium, and profit is
far off. Consider the conditions of that year and buy during the eighth
or ninth month. In a year when the harvest faces no difficulty, this is
the proper approach. In a crop-failure year, however, buying from the
sixth month onward is paramount.
\end{sokyubox}

\commentary
This chapter indicates the proper timing for speculative buying of
new-crop rice. Even if the market is at a low price, when the weather is
favorable for the season, there is no need to buy as early as the sixth
or seventh month. There may even be a good buying opportunity on a dip
caused by favorable-weather selling. However, when the weather is
unfavorable and a crop failure seems likely, one must make it a point to
begin buying from the sixth month.

% --- Article 44 ---
\article{44}{How to Manage Profits in a Rising Market}

\begin{sokyubox}
{\small\itshape\color{sokyubar} \jp{第四十四章\quad 上相塲にて利運心得方の事}}

\medskip
In a rising market, when profits are in hand and margin calls come, one
should remember the floor price and not sell back. If it falls to the
buying level of the floor, one should pile on still more. Rice that is
climbing toward a hundred-tawara rise will experience roughly two
pullbacks; at that point one must consider carefully and secure one's
profits. Moreover, when the hundred-tawara rise is near, one should
promptly close out one's position; taking profits comes first.
\end{sokyubox}

\commentary
When the market is rising and one's position has turned
profitable, and one's holdings were bought at the
floor price so that considerable profit has accrued---however tempting it
is to take profits at once, one must by no means sell here; on the
contrary, it is proper to add to the position. For example, if the
market is to rise by one full measure, two margin calls along the way
are to be expected; therefore, when the rise reaches approximately that
one full measure, one should carefully weigh the market's position and
the state of sentiment (\jt{人氣}{ninki}) and secure one's profits.
Above all, when the advance has been so great that two margin calls have
occurred, the essential thing is to take profits when profit is at hand
and then enter a new position afresh. Note that in Master S\=oky\=u's
era, a margin call was set at thirty tawara, as has already been
explained in the section on Sakata rice-exchange customs at the
beginning of this volume.

% --- Article 45 ---
\article{45}{How to Manage Profits in a Falling Market}

\begin{sokyubox}
{\small\itshape\color{sokyubar} \jp{第四十五章\quad 下相塲にて利運心得方の事}}

\medskip
In a falling market, when profits are in hand and a margin call comes,
if the counterparty elects to carry (\jt{繋ぎ}{tsunagi}) the position
rather than settle, one should promptly buy back. The first priority is
to buy back before the margin call is fully settled and the outcome
decided. This is utterly different from a rising market.
\textit{``At two, close out; at three, fully ripe; at four, it
turns''}---this is the foremost tradition of maneuvering
(\jt{駈引}{kakehiki}) on the short side.
\end{sokyubox}

\commentary
This chapter is based on experience with the trading customs of the
Sakata Rice Exchange---specifically margin calls and carrying---and
argues from that experience. Under present-day conditions, however, the
argument runs in the opposite direction. As already explained in the
section on Sakata trading customs at the beginning of this volume,
``carrying'' (\textit{tsunagi}) means that the counterparty buyer,
rather than accepting the margin call and settling, puts up additional
margin and carries the position forward. When such a carry occurs, one
should promptly buy back one's sell position and take profits. The
reason is that when a margin call is fully settled, all the sellers
scramble to buy back at once and the market rises; therefore one should
seize the initiative before the crowd takes action and take profits
first. That was the meaning of this chapter's instruction.

However, the conditions of today's market are the exact opposite, being
largely a matter of the relationship between brokers (\jt{仲買}{nakagai})
and their clients. For example, suppose the clients are all long and the
broker is short; if the broker, in order to squeeze the clients, hammers
the market down and triggers a margin call, and the clients cut their
losses and liquidate, then even in a market whose overall trend is
downward, one may see some degree of rebound. But if the clients put up
the additional margin, then since the broader trend is against them, the
market falls all the more. Such cases are plentiful, and one does well
to bear them in mind.

% --- Article 46 ---
\article{46}{New-Crop Rice Should Be Bought}

\begin{sokyubox}
{\small\itshape\color{sokyubar} \jp{第四十六章\quad 新米は買方に可付事}}

\medskip
Around the time new-crop rice first appears, one should consider each
year what price level is appropriate, according to the harvest of that
year. Autumn rice must under no circumstances be sold short. Even if
conditions do not entirely suit one's inclinations, one must without
fail take the buy side. If one truly cannot bring oneself to buy, one
should refrain from trading rather than sell. Until the annual ceiling
price has appeared, selling is absolutely useless and exceedingly
dangerous; one must take care above all not to lose the buying impulse.
A true selling opportunity exists only when the market has risen to
excess, the ceiling price is in, and there is no arbitrage rice
(\jt{乘米}{nosei-mai}) left. Apart from that, selling during a
consolidation (\jt{持合}{mochai}) market will without question result in
a loss. Rice that drifts down little by little is terrifying when it
turns upward. The reason is that because rice appears weak, both bears
and bulls sell short and open positions are numerous; therefore, when the market turns
upward, everyone scrambles to buy back at once, and the result is a
sharp, unchecked rally.
\end{sokyubox}

\commentary
This chapter is principally an admonition against selling new-crop
autumn rice.

Now, around the time each year when new-crop rice first appears, one
should consider the quality of the harvest, compare it against chart
records (\jt{罫線表}{keisen-hyo}) from prior years, and gauge the
appropriate price level. In general, autumn rice is a commodity in which
there is profit on the buy side and no profit on the sell side. The
reason is this: if the weather has been favorable that year and the
harvest is good, prices are correspondingly low, and there is no room to
sell; moreover, people talk up concerns about insufficient cutting
(\jt{鎌入不足}{kamairi-fusoku}), poor autumn weather, and the like, and
the market tends to rise. Thus the Master, too, offered the following
guidance:

\textit{``When there have been natural disasters during the summer, buy
without fail from the sixth or seventh month onward.''}

\textit{``In the sixth and seventh months, when the weather is fine,
drought continues, and there is talk of a bumper crop---when even the
slightest observation suggests a good harvest---do not rush to buy. The
eighth and ninth months are the proper time for the year's buying.''}

\textit{``Even if from the end of the sixth month through mid-seventh
month there is talk of the best harvest in decades, floods and the like
may yet occur, and the crop will suffer shortfall to some degree without
exception.''}

\textit{``Buying [autumn rice] entails months of carrying costs whose
extent is hard to know, and is indeed most disagreeable to buy; yet when
profit comes, the carrying costs are as dust before the wind and no
burden at all. Without exception, autumn rice never fails to rise.''}

And so on---he offers advice of every kind. As explained above, the
autumn market is one in which the upside prevails eight or nine times
out of ten; therefore, even if it does not quite suit one's inclination,
one should take the buy side. If one truly cannot bring oneself to buy,
one should step aside from the market entirely. Until the ceiling price
of the autumn market has appeared, selling is dangerous and must not be
attempted. One should take care to maintain a buying posture at all
times. A true selling opportunity exists only when the market has risen
to excess, the ceiling is in, and comparing prices across various
regional markets shows a relative premium. Otherwise, selling during a
consolidation market is tantamount to a guaranteed loss. Furthermore,
rice that drifts down little by little from the floor is terrifying,
because once it turns upward, it jumps to a wholly different level
without any intermediate range-trading. The reason is that because the
market looks weak, both bulls and bears have sold short and short
positions are large; once the market turns upward, everyone scrambles to
buy back at once, producing a sharp, unchecked rally.

% --- Article 47 ---
\article{47}{The Difficulty of Maneuvering on the Sell Side}

\begin{sokyubox}
{\small\itshape\color{sokyubar} \jp{第四十七章\quad 売方駈引六ケ敷事}}

\medskip
Within a single year, apart from one major rise and one major decline,
there is no great fluctuation. Within the rise there are minor declines;
within the decline there are minor advances. These are range-bound
oscillations (\jt{通相塲}{kayoi soba}) and are not markets in which one
should look for the ceiling or the floor. One must not be confused by
range-bound fluctuations. In all things the initial entry
(\jt{踏出し}{fumidashi}) is critical. Buying may cause anxiety at first,
yet when profit comes, there is not the slightest hardship. Selling, on
the other hand, feels easy at first, yet the maneuvering
(\jt{駈引}{kakehiki}) is exceedingly difficult. When prices begin to
fall, the decline seems bottomless, and one is always late to buy back.
The first priority is: do not be swept up in the downward momentum or
the crowd's panic; detach from greed and buy back.
\end{sokyubox}

\commentary
This chapter sets forth the principles of maneuvering on the sell side.
Now, within a single year---though the timing may vary---there is
generally one advance from floor to ceiling and one decline from ceiling
to floor, and no more. Within the advance there are minor dips, and
within the decline minor rallies, but these are merely range-bound
fluctuations within the larger trend---not the beginning of a true
rising market from the floor or a true falling market from the ceiling.
One must not be led astray by range-bound movement. In all things, the
initial commitment is critical: one must buy at a price close to the
floor, and sell at a price close to the ceiling. If one does not buy
near the floor, then when one of those swift corrective declines occurs
within a broader uptrend, one may panic and be shaken out of the
position. Conversely, if one has not sold near the ceiling, then when a
corrective rally occurs within a broader downtrend, one may be squeezed
and suffer terrible losses. The buyer must pay carrying charges
(\jt{利足}{riashi}) month over month to the seller, which seems a
slight disadvantage---yet when prices do rise, there is nothing to worry
about. The seller, meanwhile, receives carrying charges from the buyer,
and if profits accrue on top of that, the arrangement seems quite
favorable---yet when it comes to the seller's tactical maneuvering
(\textit{kakehiki}), the difficulty is extreme. When prices begin to
fall, sentiment turns as though the decline will never end; then a
sudden rally catches the seller off guard---because the profit-taking
buybacks of the long side come too late, the net take is diminished.
Therefore, however favorable the downward momentum may appear, and
however much the crowd insists the decline is far from over, one should
pay no heed to such clamor, suppress greed, and buy back to cover
promptly.

% --- Article 48 ---
\article{48}{Teachings Outside the Three Positions That Must Be Observed}

\begin{sokyubox}
{\small\itshape\color{sokyubar} \jp{第四十八章\quad 三位の外の教可守事}}

\medskip
The new-crop contract (\jt{新商}{shinsho})---that is, the November
contract---each year declines twenty to thirty tawara
(\jt{俵}{tawara}) from its opening price before turning upward. From
the eighth and ninth months through the tenth, eleventh, twelfth, and
first months, it invariably rises. However, in a year when Ky\=ush\=u
and the Kamigata region produce a crop around eighty percent or better,
no such rise occurs. When natural disasters strike, prices surge sharply
from the seventh or eighth month onward. When the harvest shortfall is
gradual and natural, the nationwide average crop yield becomes known in
\=Osaka by the ninth month, and from the tenth through twelfth months
\=Osaka prices are pulled upward, so local prices rise as well. Consider
this carefully.

In years when both the Kamigata and local harvests are average, or when
the local harvest is good, there is only a modest rise through autumn
and winter---no great advance. In a year of local crop failure,
regardless of conditions in the Kamigata or elsewhere, local prices
surge dramatically. In years of four-and-a-half- to six-tenths crop
failure, one should anticipate that prices will eventually reach the
ten-tawara level, and reckon carefully when they approach twenty tawara.
The ceiling price, once reached, does not fall for fourteen or fifteen
days, then declines gradually of its own accord. Through the third and
fourth months, prices oscillate within a range of about ten tawara,
declining twenty to thirty tawara; then a consolidation (\jt{保合}{mochai})
follows. By the fourth and fifth months prices fall thirty to forty
tawara further, then rally twenty to thirty tawara again. Come the sixth
month, with fair weather, splendid reports from both paddies and dry
fields, both domain-warehouse rice (\jt{御藏米}{okura-mai}) and private
stocks remaining abundant, and prices too high for the shipping buyers
(\jp{船々}, \textit{funabune}) to rush in---at this point those holding
long positions grow weary and sell out, and prices suddenly collapse
sixty, seventy, even a hundred tawara. One must understand that this is
the great crash.
\end{sokyubox}

\begin{sokyubox}
{\small\itshape\color{sokyubar} Article 48 (continued)}

\medskip
Furthermore, in a year when the fifth month is the head month of the
Three Positions (\jt{三位}{sanmi}), if prices do not decline
sufficiently in the fourth and fifth months, there will be no rise in
the sixth or seventh months. When prices do decline sufficiently across
the fourth, fifth, and sixth months, prices will certainly rise in the
seventh and eighth months---there is no doubt of it. When the seventh
month is the head month of the Three Positions, and prices have not
risen in the fourth and fifth months, the sixth month is the floor
month, and prices will certainly suffer a great crash by around the
twentieth. When one has sold from early in the sixth month and prices
begin falling, consider the floor to lie in the interval around the
twelfth or thirteenth day; if margin calls (\jt{落引}{ochibiki}) are
approaching, that is one thing, but if prices have not yet fallen,
identify the floor, buy back, and go further to establish fresh long
positions. Then wait for the subsequent rise, and when profit appears,
close out promptly.

Rice that makes its ceiling price in the seventh, eighth, ninth, or
tenth months will decline by the twelfth month---observe this rule. Rice
that makes its floor price in those same months will rise by the twelfth
month---observe this teaching. Rice that makes its ceiling price from
midwinter through the first or second months will decline by the fourth,
fifth, or sixth months. When prices fall sufficiently in the fifth
month, the sixth month brings a sharp rally. When prices do not fall
sufficiently in the fifth month, there will be a great crash by around
the twentieth of the sixth month. In years when prices have been at low
levels since the previous winter, speculative long positions pile up,
leading to a sixth-month crash. Rice that has been at depressed
levels from midwinter
through the second or third months will certainly rally in the
summer---one should expect it to rise as late as the seventh month.
These are the teachings outside the Three-Position Tradition
(\textit{sanmi no den}); take them to heart. In any case, autumn rice
should not be sold short---one should take the long side. Regardless of
the magnitude, it invariably rises. As for the falling market after a
ceiling price, that matter is recorded in the section on the Three
Positions; study it well and put it into practice.
\end{sokyubox}

\commentary
This chapter is titled ``Teachings Outside the Three Positions That Must
Be Observed,'' but when one reads it carefully and digests its meaning,
it is a concise synthesis of the market principles expounded in the
preceding chapters---something like a verse (\textit{g\=ath\=a})
appended to a Buddhist s\=utra.

Considering the normal pattern: the new-crop contract (i.e., the
November contract) opens each year around the sixth or seventh month.
Customarily, prices decline twenty to thirty tawara from the opening
price, and then turn upward. From the eighth and ninth months through
the tenth, eleventh, twelfth, and first months, they seem never to fail
to rise. However, in a year when the Ky\=ush\=u and Kamigata crops are
around eighty percent or better, no such vigorous rise occurs---the
market is weak and listless. But when natural disasters strike, prices
surge from around the seventh or eighth month. In years when cold
weather and unseasonable conditions produce a naturally poor harvest, as
the nationwide crop conditions become known, T\=oky\=o prices rise from
the tenth through the twelfth month, and each locality rises in
response.

In years when both the Kamigata and the local area produce an average or
above-average crop, there is only a slight rise through autumn and
winter and no great advance. In years of local crop failure, regardless
of average or above-average harvests in the Kamigata or elsewhere, local
prices alone surge dramatically. In years of four-and-a-half- to
six-tenths crop failure, one should look ahead to prices reaching the
ten-tawara level and make careful judgments when prices approach twenty
tawara. In such years of severe crop failure, even when the ceiling
price is reached, prices do not drop sharply as they usually would, but
hold at high levels for fourteen or fifteen days in consolidation, then
decline gradually.

The typical pattern is: through the third and fourth months, prices
oscillate within a range of about ten tawara, fall twenty to thirty
tawara, and enter a consolidation; by the fourth and fifth months they
decline thirty to forty tawara, then rally twenty to thirty tawara
again; then in the sixth month, with fair weather, both the green rice
paddies and the wheat crop looking splendid and reports of a bumper
harvest, both domain-warehouse rice and farm-yard rice (\jt{庭米}{niwa-mai})
in ample supply, and the price too high for shipping buyers to rush
in---seeing all this, those who had bought on speculation and those
holding physical rice grow weary and begin selling out, whereupon prices
suddenly collapse sixty to seventy tawara or even close to a hundred
tawara.

Furthermore, in a year when the fifth month is the head month---that
is, the first-stem month (\jt{甲月}{kinoe-tsuki})---if prices do not
fall sufficiently during the fourth and fifth months, there will be no
rise in the sixth or seventh months. Conversely, if prices do fall
sufficiently across the three months from the fourth through the sixth,
the seventh and eighth months will without doubt bring a rise. When the
seventh month falls on the first-stem month, and prices have not
declined in the fourth and fifth months, the sixth month has come to
function virtually as a perennial floor month; prices will suffer a
great crash by around the twentieth.

The tactical response is as follows: sell from early in the sixth month,
and once prices begin falling, watch the period from around the tenth to
the twelfth or thirteenth day to gauge where the floor lies. Take profit
promptly by buying back, and go further to establish fresh long
positions; then wait for the subsequent rise and, again, take profit
promptly by selling out.

The rules to keep in mind are these: rice that makes its ceiling price
in the seventh, eighth, ninth, or tenth months will decline by the
twelfth month. Rice that makes its floor price in those months will rise
by around the twelfth month. Rice that makes its ceiling price from
midwinter through the first or second months will decline by the fourth,
fifth, or sixth months. When prices fall sufficiently in the fifth
month, the sixth month brings a sharp rally; when they do not fall
sufficiently, prices collapse by around the twentieth of the sixth
month. Rice that has been at depressed levels from the previous winter
through the second or third months is certain to rally in the
summer---sometimes continuing as late as the seventh month.

These are the teachings outside the Three-Position Tradition explained
in the preceding and following chapters. In any event, autumn rice
offers no profit on the short side; regardless of the degree, it
invariably rises, so one should take the long side. These explanations
have already been given in the preceding chapters and will not be
repeated at length. As for the falling market following a ceiling price,
that too has been explained in detail in the section on the Three
Positions; refer to it there.

% --- Article 49 ---
\article{49}{When Summer Brings Cold Weather, One Must Not Let Down One's Guard}

\begin{sokyubox}
{\small\itshape\color{sokyubar} \jp{第四十九章\quad 夏中冷数時は油断ならぬ事}}

\medskip
In a summer of bad weather and persistent rain, one should understand
this as a natural crop failure. When summer cold is severe and there are
insect infestations, floods, great storms, and other natural disasters,
one must not let down one's guard---buy into the market while prices are
still low from the sixth and seventh months. In the sixth and seventh
months, when the weather is fair and drought persists, when reports
speak of a bumper harvest and every scrap of information points to a
good crop, do not rush to buy. The eighth and ninth months are primarily
the year's buying opportunity. Buy at the opening of each month. The
ninth month tends to be slightly late; focus principally on the eighth
month. Even in a bumper-crop year, prices invariably rise from the ninth
and tenth months onward. Furthermore, even if from late in the sixth
month through mid-seventh month the talk is of the best crop in decades,
something or other will emerge, and the harvest will fall short to some
degree without fail. A harvest that comes in without the slightest
shortfall is exceedingly rare---consider this well.
\end{sokyubox}

\commentary
This chapter concerns the timing of speculative buying of the new-crop
contract. The meaning is as follows: when the weather is unseasonable
and the prospect of a poor harvest is well established, it is proper to
begin buying from the sixth or seventh month. But in a year of fair
weather, sentiment-driven declines may appear from mid-seventh month
into the eighth month, so one should not rush to buy. Rather, target the
very bottom of the listless drift---the buying opportunity. Because
supplementary settlement discounts (\jt{補安}{hoyasu}) tend to bring
prices down at the opening of each month, those are the levels to buy.
The buying opportunity appears most often in the eighth month.

Even in any bumper-crop year, prices invariably rise from around the
ninth or tenth month. Even if from late in the sixth month through
around mid-seventh month people proclaim the best harvest in decades,
afterwards some form of natural disaster appears, the harvest falls
short to some degree, and prices rise accordingly. A crop that reaches
full harvest without the slightest hitch is exceedingly rare---weigh
this well.

% --- Article 50 ---
\article{50}{Guidance for Bumper-Crop Years Across All Regions}

\begin{sokyubox}
{\small\itshape\color{sokyubar} \jp{第五十章\quad 諸方上作年心得の事}}

\medskip
When the local region and neighboring provinces have a bumper crop, and
the Kamigata likewise produces largely a bumper crop, yet no floor price
(\jt{底直段}{soko nedan}) has emerged by the third, fourth, or fifth
months---one should understand that the floor price will emerge
thereafter.
\end{sokyubox}

\commentary
In years when the local district and the Mogami region both enjoy an
excellent harvest, there is a sell-off around the middle days of the
fourth month---this is beyond doubt. Even if the third month falls on
the Three Positions, there was a crash of about twenty tawara in the
fourth month of Ky\=owa~2 (\jp{享和二年}, 1802, Year of the Rooster).
In years when nearby provinces and the local area have a bumper crop,
one must absolutely not establish speculative positions in summer rice
until the first weeding (\jp{一番草取}, \textit{ichiban kusatori}) is
past. One should, however, be prepared to buy on the fourth- and
fifth-month crashes.

When both the local area and neighboring provinces have a bumper crop,
even if a considerable decline occurs, various circumstances may prevent
the market from falling through---prices remain in a relatively high
consolidation without producing a proper floor. Yet the market does
produce a crash and a floor-like bottom during the three months from the
third through the fifth. The case of Ky\=owa~2 is one such example.

In bumper-crop years for both the local area and nearby provinces, one
must not establish speculative summer-rice positions until the first
weeding season has passed. Focus on buying the fourth- and fifth-month
crashes.

% --- Article 51 ---
\article{51}{The Seventh Month as the Leading Month of the Three Positions}

\begin{sokyubox}
{\small\itshape\color{sokyubar} \jp{第五十一章\quad 七月三位頭月の事}}

\medskip
When the seventh month falls on the leading month (\jt{頭月}{kashira-tsuki})
of the Three Positions, even though the sixth month may have fine weather
and everything points to an excellent crop with further room to fall,
pay no heed to that---if rice is around thirty tawara, buy without
hesitation from the sixth month onward. As for autumn rice, even when
all Japan proclaims a bumper harvest, it is rare that the crop escapes
difficulty entirely. Should it be an average year with both the local
area and the Kamigata enjoying good harvests, then hold off from the
eighth, ninth, and tenth months, identify the floor price, and buy in.
In most years of ordinary harvests, prices decline ten to twenty or
thirty tawara from the opening and then mostly turn upward. A decline of
forty to fifty tawara is exceedingly rare. However, if during the sixth
and seventh months unexpected natural disasters strike both the Kamigata
and the local area, there will be no orderly trading---prices leap
upward by the bale in a sudden surge, and one is left behind. One must
absolutely not let down one's guard. This matter is of the greatest
importance; engrave it upon the mind. In short, when the seventh month
is the leading month, buying from the sixth month onward is without
error eight or nine times out of ten. The seventh month as leading month
is a month when unusual events frequently occur in both the Kamigata
and the local area---it is a month of sharp rallies. One must never be
complacent.
\end{sokyubox}

\begin{sokyubox}
{\small\itshape\color{sokyubar} Article 51 (continued)}

\medskip
When the eighth month falls on the ceiling and terminal month
(\jt{終月}{owari-tsuki}), it is rare for either the new-crop or
old-crop contract to produce a ceiling in this month. The old-crop
contract does end in this month, to be sure. But the new-crop contract
fluctuates through the seventh month of the following year, so one
cannot fix the eighth month as its ceiling. This too tends to fall in
years when the Ky\=ush\=u region is due for natural disasters. Even if
the eighth month does produce a ceiling price, in a year when Ky\=ush\=u
and the Kamigata suffer crop damage, there will be a rise approaching a
hundred tawara within the tenth or eleventh month (\jt{霜月}{shimotsuki})---careful
judgment is paramount.

Rice that has struck a ceiling at around eighteen to twenty-one or
twenty-two tawara, and by the following summer has declined only forty to
fifty tawara before buyers come in, indicates that physical rice in the
Kamigata is in short supply and shipping buyers are rushing to purchase.
In the normal course, one would expect a decline of sixty or seventy to
nearly a hundred tawara from the first-stem month through to the floor
month of the Three Positions. Yet in a year when sentiment remains
stretched from the prior year's poor harvest and prices fail to decline
as far as they should through the fourth and fifth months, once the busy
shipping season passes its peak, prices ease slightly; in the fifth
month, when shipping pauses, they fall twenty to thirty tawara further.
Then in the sixth month, with fine weather and strong sunshine through
the midsummer \textit{doy\=o} period, reports of a bumper crop emerge;
prices are too high relative to the cost basis, so shipping buyers hold
back; sentiment deteriorates steadily. The sixth month is also the month
when rice holders close out their positions---outside speculators with
long positions, as well as non-professional and back-channel traders
(\jp{素人黒人}, \textit{shir\=oto kur\=oto}), all sell to close. From
mid-sixth month onward a great crash is beyond doubt. Especially in a
year when the previous year saw a four-and-a-half-tenths crop failure
and outside speculative longs are heavily accumulated, the collapse is
all the greater. One should consider the circumstances of such years
very carefully.

That said, even in a year when the eighth month corresponds to the
ceiling, one should not chase high prices. Wait for a floor-like dip
during the three months from the fourth through the sixth, buy both
old-crop and new-crop rice, and wait for the eighth-month rally to take
profit.

In most cases, an eighth-month ceiling is rare---it does not occur in
ordinary years. When a ceiling price appears in winter or spring followed
by an extraordinarily steep decline, another ceiling will emerge in the
seventh or eighth month. The Tradition says, \textit{``When prices fall
sufficiently in the fifth month, the sixth month brings a sharp
rally''}---the logic is the same. In an ordinary year, when a summer
rally runs from the third and fourth months through the fifth and sixth
months, one must understand that there will certainly be no rise in the
eighth month.
\end{sokyubox}

\commentary
This chapter explains the principles to follow when the seventh month
falls on the first-stem month (\textit{kinoe-tsuki}). The meaning is as
follows: when the seventh month is the first-stem month, even if the
sixth month has fine weather and the outlook points to a bumper crop with
apparent room for prices to fall further, one should pay no attention to
such appearances. If prices are around thirty tawara, buy the dips
without hesitation. Even if the whole of Japan proclaims a bumper
harvest during the summer, nine years out of ten, come autumn, complaints
of crop difficulty emerge and prices rise.

Should the previous year also have been an average year with good
harvests in both the local area and the Kamigata, one need not rush to
buy but may wait until the eighth, ninth, or tenth month, identify the
floor price, and buy in then.

Having determined the floor, one should buy in. In ordinary years, when
the crop is about average, from the opening price of the new contract
onward, there is typically a decline of ten to twenty or thirty tawara,
after which the market mostly turns upward. A decline of forty or fifty
tawara from the opening price is exceedingly rare. However, during the
sixth and seventh months, when an unexpected natural disaster strikes
both the Kamigata region and the local area, the market breaks sharply
upward without any conspicuous trading---prices jump by the bale and
surge higher---and one is left behind. One must never be caught off
guard. This point is of the utmost importance and must not be forgotten.

In short, experience shows that when the seventh month falls on a
first-stem month (\jp{甲}, \textit{kinoe}), buying from the sixth month
onward proves unprofitable only one time in ten. When the seventh month
falls on a \textit{kinoe} month, it is generally a month in which some
disturbance arises in both the Kamigata and the local area, producing a
sharp rally---one must never let one's guard down.

Next, regarding an eighth-month ceiling: it is truly rare for the
ceiling to appear in this month. Since the Sakata Rice Exchange's
(\jt{酒田米會所}{Sakata Kome Kaisho}) customary trading season for old
rice ends on the last day of the eighth month, and from the ninth month
onward only new-rice contracts are traded, and since new rice fluctuates
through the following seventh month, there is no particular reason, on
principle, why the ceiling must appear in this month. A \textit{kinoe}
seventh month tends to coincide with a cycle in which natural disasters
strike the Ky\=ush\=u region. Even if a ceiling price appears in the
eighth month and the market eases temporarily, should Ky\=ush\=u and
the Kamigata suffer crop damage that year, there will likely be a rise
of around one hundred tawara by the tenth or eleventh month---one should
watch the situation carefully.

When rice that has struck a ceiling at eighteen to twenty-one or
twenty-two tawara shows, the following summer, only a modest decline of
forty or fifty tawara, and buying ships are pressing in, this is because
physical rice in the Kamigata and the provinces is in short supply, and
the ships are rushing to buy. Properly speaking, a decline of sixty or
seventy to roughly one hundred tawara from the first-stem month through
the last-stem month would be the normal course. Yet in years when strong
sentiment persists from the previous year's poor harvest and prices fail
to decline as far as they should by the fourth or fifth month, from the
busy shipping season onward the market turns slightly lower; in the
fifth month, when shipping briefly thins out and there are no buyers, it
falls another twenty or thirty tawara. Then, in the sixth month, when
the weather is fair and summer heat beats down during the \textit{doy\=o}
period, good crop reports circulate; moreover, because prices are
disproportionately high and the arithmetic does not work for buyers in
the provinces, the ships do not buy, and sentiment gradually
deteriorates. The sixth month is also the month when rice holders sell
off their inventory---not only provincial speculators but amateurs and
professionals alike close out---and from mid-sixth-month onward a
collapse is beyond doubt. Especially in years following a poor harvest,
when provincial speculative positions are heavily built up, the collapse
is all the greater. In such years, one should pay very close attention.

However, even in a year when the eighth month falls on a \textit{kinoe}
ceiling month, it does not pay to chase high prices. During the three
months of the fourth, fifth, and sixth, whether in new rice or old, one
should look for a floor-like dip and buy in, then wait for the
eighth-month rally to take profit (\jt{利喰}{rigui}).

That said, judging by past precedent, an eighth-month ceiling is
exceedingly rare---it simply does not occur in ordinary years. Only in
years when the ceiling price has appeared from winter through spring and
the market has subsequently fallen to an extraordinary degree does the
ceiling reappear in the seventh or eighth month.

Beyond that, the kinds of years in which a seventh- or eighth-month
ceiling may arise are those in which, around the sixth month, the
weather is unfavorable and there is anxiety about a poor harvest ahead.
In such years, by the custom of the Sakata Rice Exchange, the opening of
the new-rice contract is postponed from the sixth month. If the market
stands at the thirty-tawara level, one should buy during the fifth and
sixth months while prices are low. For physical rice
(\textit{sh\=omai}), there are risks of spoilage and measure-shrinkage,
as well as interest due to the selling counterparty; estimate these
costs at twenty to thirty tawara and use rice certificates as collateral
for borrowing (\jt{歩金}{bukin}---pledging rice certificates as
security, much like pawning), then wait for the eighth-month rise.
Should an extraordinarily high price appear in the seventh month, one
need not wait until the eighth---take profit then. This is the same
principle described earlier: ``When the fifth month falls fully, the
sixth month rallies sharply.'' In ordinary years, when there is a summer
rally from the third or fourth month through the fifth or sixth, there
will decidedly be no eighth-month rise---one should understand this
well.

% --- Article 52 ---
\article{52}{On the Great Collapse from Winter through Spring When Cash Is Depleted}

\begin{sokyubox}
{\small\itshape\color{sokyubar} \jp{第五十二章\quad 冬より春迄金拂底大崩しの事}}

\medskip
In a year when cash is depleted (\jt{金拂底}{kinbarai soko}) from
winter through spring, it is because speculative buying is heavy and
funds run short. If in such a year prices have not fallen by the fifth
month, a great collapse around the twentieth of the sixth month is
beyond doubt. Should the seventh month in such a year fall on a
\textit{kinoe} month, prices should rise sixty or seventy tawara in the
seventh month.

In Meiwa~5 (\jp{明和五年}, 1768), the decline from the ceiling price
was one hundred sixty to one hundred seventy tawara---a decline greater
than this is rare. In Kansei~11 (\jp{寛政十一年}, 1799), it was one
hundred fifty tawara down; in Kansei~12, on the thirteenth day of the
sixth month of the Monkey year, one hundred tawara down. The floor price
persisted for three years.

When the seventh month falls on a \textit{kinoe} month and the market
turns downward, sell without hesitation; watch for twelve or thirteen
days, target a hundred-tawara decline, then buy back and buy beyond that.
\end{sokyubox}

\commentary
In Master S\=oky\=u's era, financial institutions were not as developed
as they are today, and the avenues for deploying capital were nothing
like what they are now. Consequently, in his domain and the local area,
the uses of capital were limited to speculating in rice or speculating in
other commodities. Considered from this standpoint, in a year when
currency in circulation is depleted from winter through spring, it is
proof that speculative buying in rice is heavy. In such a year, if the
market has not declined by around the fifth month, it will assuredly
collapse by late in the sixth month. Should the seventh month that year
fall on a \textit{kinoe} month, prices should rise sixty or seventy
tawara in the seventh month.

``When the market turns downward, sell without hesitation. The first
ceiling is difficult to identify, so watch for twelve or thirteen days,
then sell without fail at the second ceiling. Target a hundred-tawara
decline, assess the sentiment at that point, take profit promptly, and
at times buy beyond the close.''

% --- Article 53 ---
\article{53}{On Ceiling and Floor Prices Persisting for Three Years}

\begin{sokyubox}
{\small\itshape\color{sokyubar} \jp{第五十三章\quad 天井直段底直段三年續きの事}}

\medskip
Ceiling prices and floor prices, it is said, persist for three years.
One must not suppose this applies only to ceiling prices---floor prices,
too, persist for three years. Since this occurs even in ordinary years,
one should pay very close attention.
\end{sokyubox}

\commentary
In the \textit{Three-Position Catechism} of the Kuzuoka-gura
(\jp{葛岡倉の三位間答}, \textit{Kuzuoka-gura no sanmi mond\=o}), this
is explained as follows: if the market cycles high for three years, then
the low cycle likewise continues for three years. Between these, there
are middling years that belong to neither the high cycle nor the low.
Cycling in three-year intervals, approximately once every ten years a
great high cycle occurs.

This accords with the theory current in Europe and America that a panic
occurs once every ten years. More detailed discussion will follow later,
but in the passage elsewhere in this text concerning ceiling prices, it
says: ``In the first year, a ceiling of one hundred to one hundred
twenty, one hundred sixty, or one hundred eighty tawara may appear; in
the second year, it falls short of the first; in the third year, it
falls still shorter---understand this well. The pattern in the following
spring's decline is the same.'' One should consult that passage as a
reference.

% --- Article 54 ---
\article{54}{On Gauging the Ceiling: Targeting a Hundred-Tawara Rise}

\begin{sokyubox}
{\small\itshape\color{sokyubar} \jp{第五十四章\quad 天井の考百俵上げ的の事}}

\medskip
To gauge the ceiling means targeting a rise of one hundred tawara or
more. To gauge the floor means targeting a hundred-tawara decline---yet
one may be confounded by the prevailing conditions. When prices begin to
fall, the decline continues for ten, twelve, or thirteen days; therefore,
watch carefully and wait for the bottom before buying in. On \textit{ushi}
days (\jp{丑の日}), \textit{fuj\=oju} days (\jt{不成日}{fuj\=oju nichi}---inauspicious
completion days), \textit{ten'ichi} days (\jp{天一}), and \textit{hassen}
days (\jp{八專}), the decline most often halts. In
particular, when buying in a falling market, unless one buys the floor,
profits will be thin.
\end{sokyubox}

\commentary
That the ceiling and floor should be determined by a hundred-tawara
swing as the standard has been stated repeatedly. This was based on the
commercial customs of the Sh\=onai region in S\=oky\=u's time. In the
present day, following the standard established by Master Ishikawa, one
should treat roughly one hundred tawara as approximately one-tenth and
conduct one's maneuvering (\textit{kakehiki}) accordingly.

Now, to gauge the ceiling and floor, one targets a swing of one-tenth or
more. There are times when one is swept up in the sentiment of the
moment and finds it difficult to sell into the ceiling---this is not
uncommon. But regardless of the frenzy of the moment's sentiment, when
one detects the first sprout of a decline, one should sell without
hesitation. Furthermore, when rice begins to fall from the ceiling, the
decline continues for ten to twelve or thirteen days in succession. To
scoop up small lots, one should wait for each successive dip. There are
days when the decline pauses of its own accord: on \textit{ushi} days,
\textit{fuj\=oju} days, \textit{ten'ichi} days, and \textit{hassen}
days, the decline most often halts.

One should consult the passage elsewhere in this text that says:
``\textit{Ushi} days and \textit{tatsu} days are days when rice is
cheap; however, \textit{ushi} days are days when floor prices and
ceiling prices appear---at ceilings and floors, pay close attention to
\textit{ushi} days.'' Next, regarding the statement that ``in a falling
market, unless one buys the floor, profits will be thin''---in a falling
market, scooping up small lots on the way down is difficult to profit
from, because while one is still thinking the retracement has not gone
far enough, the market turns and heads down again.

% --- Article 55 ---
\article{55}{On Reading the Sixth-Month Collapse}

\begin{sokyubox}
{\small\itshape\color{sokyubar} \jp{第五十五章\quad 六月崩し見様の事}}

\medskip
When the previous year's harvest was poor, or when both the Kamigata and
the local area suffered poor harvests, and there are rumors of goods in
short supply, and \=Osaka prices are high---at such times, local
sentiment is fired up, speculative buying of summer rice is heavy, and
cash runs short. When buying is heavy, take heed. Around the eleventh
month, twelfth month, or first month, a ceiling of eighteen or nineteen
to twenty-two or twenty-three tawara appears. Rice that then holds in
consolidation (\jt{保合}{mochai}) through the fifth or sixth month will
assuredly collapse in the sixth month.

When a ceiling appears by around the tenth or eleventh month and there
has been a rise of sixty or seventy tawara by the twelfth month or first
month---and when the previous year was a half-crop or worse and a great
ceiling of seventeen or eighteen tawara appears in winter---sentiment is
stretched to the extreme. Because summer-rice speculation is intense,
prices do not decline through the fourth or fifth month. In the sixth
month, the weather is fine, the sun beats down through the \textit{doy\=o}
period, and reports of splendid paddy growth circulate. Since prices are
generally high, the buying ships do not press; meanwhile, because both
domain warehouse rice (\jp{御藏}, \textit{okura}) and private rice
(\jt{内米}{uchimai}) come to market in insufficient quantities, these
rumors cause sentiment to deteriorate, and sellers rush to escape, each
trying to be first. Margin deposits (\jt{元質}{motoichi}) also fall
short, and demands for additional funds, combined with the rest, bring
abundant supply onto the market---and prices fall further still.

One must first of all consider the volume of ships, the volume of rice
coming to market, and the size of long positions. \textit{``What falls
short will prove to be in excess; what is in excess will prove to fall
short''}---in time, when the supply has thinned out, prices will rise.
Consider this well.

However, when one judges that a sixth-month collapse is likely, one
should assess conditions carefully and sell from the start of the sixth
month through around the tenth day, even at a thirty-tawara decline.
Then observe the pace of the decline for six or seven to twelve or
thirteen days from the start, gauge the floor price, and buy in. Because
the season is late, large positions are pointless---limit oneself to two
or three hundred ry\=o. In a year when the seventh month falls on a
\textit{kinoe} month, one may be confident. Otherwise, years in which
buying ships ply their trade through the eighth month are rare---judge
accordingly.
\end{sokyubox}

\commentary
This chapter describes the conditions and the proper approach to a
sixth-month collapse.

When the previous year's harvest was poor, and conditions in the
Kamigata likewise lean toward poor harvests with rumors of physical rice
in short supply, and the \=Osaka market is at high prices---local
sentiment is stimulated and speculative buying of summer rice is heavy.
The depletion of cash is the proof. At such times, one must pay the
closest attention. When a ceiling of eighteen or nineteen to twenty-three
or twenty-four tawara appears around the eleventh month, twelfth month,
or first month, and rice holds at high-price consolidation through the
fifth or sixth month, a sixth-month collapse is certain.

Again, when a first ceiling is struck by around the tenth or eleventh
month, and thereafter the ceiling falls sixty or seventy tawara, and the
market holds in consolidation around the twenty-seven-tawara level
through around the fourth month---such rice will rally sixty or seventy
tawara in the seventh or eighth month.

In a year when the previous year's harvest was roughly half-crop and a
great ceiling of seventeen or eighteen tawara appears, sentiment is
extremely taut, and because summer-rice speculation is heavy, prices do
not decline through the fourth or fifth month. From around the
\textit{doy\=o} period of the sixth month, the weather is praised and
buyers hold back anticipating a good harvest, while neither the domain's
warehouse rice nor private supplies come out in sufficient quantity. As
the market declines step by step, margin deposits on physical rice used
as loan collateral also fall short, and the disposition of that rice
comes onto the market too---all told, supply overwhelms demand and
prices fall relentlessly.

At such a time, one must consider the volume of buying ships, the volume
of rice coming to market, and the scale of long positions, and maneuver
accordingly. Yet as the proverb says, \textit{``what falls short will
prove to be in excess; what is in excess will prove to fall
short''}---in time, physical rice thins out and prices rise. One should
consider this well.

When one judges that a sixth-month collapse is developing, one should
sell from mid-fifth-month onward, even if prices seem somewhat low, and
when the collapse begins, gauge the floor and take profit within about
half a month. However, because the season is late, one must not commit
large capital.

% --- Article 56 ---
\article{56}{On Resolve: ``Just as I Thought'' Whether It Rises or Falls}

\begin{sokyubox}
{\small\itshape\color{sokyubar} \jp{第五十六章\quad 上れば扨こそ下れば扨こそ决心の事}}

\medskip
One may believe that rice should rise, yet when an unexpected dip
appears in the interim and one watches uneasily, and then prices rise a
little, one thinks ``just as I thought---it is going up.'' When prices
then fall, one thinks ``just as I thought---it was going down after
all.'' This is nothing but a mind that is unsettled and in constant
motion.

When one studies the market in hindsight, it seems easy enough: ``I
could have sold that rally and bought that dip, capturing both the rise
and the fall.'' But in practice, capturing both sides of the swing is
rare. Rather, one should step back for a month or two, observe the
aftermath, and resolve to capture just one movement---either the rise or
the fall. The desire to capture both sides of every swing keeps one
ceaselessly tied to the market the whole year through, one's mind in
constant agitation. That is why losses mount.
\end{sokyubox}

\commentary
In trading, one must empty one's mind and make it level, then discern
whether prices are headed up or down---this is what matters most. When
one's mind is even slightly unsettled, the merest uptick makes one think
``Ah, so it is a rising market after all,'' and the merest dip makes one
think ``Ah, so it was a falling market''---and one's mind knows not a
moment's rest. Looking at the market's footprints after the fact and
thinking about them, it seems as though one could have profited on both
the rise and the fall. But that is the advantage of the spectator---like
an onlooker at a game of go who sees eight moves ahead. When one is
actually in the thick of it, far from capturing both sides, one cannot
even capture one.

Therefore, abandon such greed. Rather than chasing the market's every
move, step back for a month or two, form a longer view, and resolve to
capture just one direction---up or down. As stated earlier, the
principle of committing to one side---buy only when you see an advance,
sell only when you see a decline---applies precisely here.

This chapter pierces the weak point of the market speculator with
penetrating insight. The reader will surely recognize having found
himself in just such a predicament more than once.

% --- Article 57 ---
\article{57}{On the Error of Averaging Down in a Bearish Losing Streak}

\begin{sokyubox}
{\small\itshape\color{sokyubar} \jp{第五十七章\quad 弱氣不利運心得違の事}}

\medskip
When one takes a bearish stance and turns to the sell side, and through
misjudgment begins to accumulate small losses, one may then attempt to
sell-average (\jt{賣平均}{uri-beikin}) by piling on more sell orders
into rice that is already rising. This is a grave error of judgment. One
must not go against the market; exercise restraint. The same applies
when on the buy side. When one's assessment proves wrong, close out
early and observe how things develop.
\end{sokyubox}

\commentary
This chapter is a warning against averaging into a losing position---sell-averaging
or buy-averaging, that is, the practice of adding to a losing trade
(\jt{難平商}{nanpin akinai}). The meaning is clear without further
annotation.

% --- Article 58 ---
\article{58}{In a Bountiful Year, Rice Must Not Be Sold}

\begin{sokyubox}
{\small\itshape\color{sokyubar} \jp{第五十八章\quad 豐年の米不可賣事}}

\medskip
When the harvest in this region is bountiful, one must not sell rice.
During autumn and winter, when the floor price (\jt{底直段}{soko nedan})
appears, buy in and carry the position through to spring; profit is
assured. Because the news of a bountiful harvest reaches the other
provinces, by spring incoming shipments (\jt{入船}{irifune}) increase,
and winter buying also arrives. During autumn and winter, prices are
pulled down in step with the crop yield. If prices do not rise by
spring, they will certainly rise by the fifth or sixth month.
\end{sokyubox}

\commentary
In a year of bountiful harvest in this region, the market will already
have discounted the surplus, driving prices down close to the
floor---so taking the sell side is unprofitable. In
a year of good harvest, the floor price will inevitably appear sometime
during that autumn or winter; one should buy at that point and carry the
position to the following spring, when profit is assured. The reason is
that when word of cheap prices from a good harvest spreads to the other
provinces, winter buying flows in and, moreover, by spring the incoming
shipments increase, creating the conditions for a spring rally. In some
cases, depending on circumstances, the spring rally may be delayed until
the fifth or sixth month. For further detail on incoming shipments and
winter buying, see the section on Sakata rice-trading customs at the
beginning of this volume.

% --- Article 59 ---
\article{59}{On the Two Kinds of Regret}

\begin{sokyubox}
{\small\itshape\color{sokyubar} \jp{第五十九章\quad 後悔に二つある事}}

\medskip
There are two kinds of regret in trading; one must understand both.
When prices are fluctuating and one could have taken full profit by
waiting five or six more days, but rushes to take profits and captures
only two- or three-tenths---this is the regret one laughs off. On the
other hand, when one holds rice showing seven- or eight-tenths profit
but, lost in greed, cannot bring oneself to close, and the market
reverses into a loss---this is the regret that comes after great toil,
the kind that truly wounds the spirit. One must take this to heart and
be on guard.
\end{sokyubox}

\commentary
In the matter of profit-taking (\jt{利喰}{rigui}) on the market, there
are two distinct categories of regret. One is the regret one laughs off;
the other is the regret that wounds the spirit. Examples of each are
detailed in the main text. However, the tactical play of the market
demands measures adapted to the moment, and one cannot lay down a fixed
rule declaring the former approach always better or the latter always
preferable. If one always takes only a small portion of the profit, one
will never achieve a large gain; yet if one always holds out greedily,
trying to capture twelve-tenths of the profit, sooner or later one will
be caught by a reversal and lose the gains already in hand. One must
carefully read the broad trend of the market and the waxing and waning
of sentiment (\jt{人氣}{ninki}), and take measures as the moment
requires. In sum, the chapter's meaning comes down to this: when one
holds a profitable trade, close out at seven- or eight-tenths.

% --- Article 60 ---
\article{60}{On the Need for Daily Attention to the Market}

\begin{sokyubox}
{\small\itshape\color{sokyubar} \jp{第六十章\quad 毎日の相塲に氣を付可申事}}

\medskip
The ultimate secret is to pay attention to the market every day and to
trade according to the oscillation of rice. What is referred to here as
``oscillation and advance'' (\jt{通ひ運び}{kayoi hakobi}) is a secret
tradition. In reading this, the difference between right and wrong is
stark---a hair's breadth of error leads a thousand leagues astray. This
is the first principle.
\end{sokyubox}

\commentary
What the text calls ``oscillation'' (\jp{通ひ}, \textit{kayoi}) refers
to the range-bound market (\jt{通相塲}{kayoi soba}), and ``advance''
(\jp{運び}, \textit{hakobi}) refers to the trending market
(\jt{伸相塲}{nobi soba}).

Now, the ultimate secret of the rice trader is to pay daily attention to
the market and judge whether the current market is one that oscillates
back and forth in a holding pattern (\jt{保合}{mochai}), or whether it
is trending upward or crashing downward as a trending market. One should
attempt trades based on this judgment. However, the distinction between a
range-bound and a trending market is extremely difficult to make; even a
slight error in reading can have consequences for one's entire fortune,
so one must exercise the utmost care in this observation. Let us here
synthesize the master's teachings and say a few words about both types of
market.

A trending market means a major market that climbs steadily from the
floor price or crashes from the ceiling (\jt{天井}{tenj\=o}). Its
standard is the monthly high-low mark (\jt{月節}{tsukibushi}): a market
that breaks above or below the prior month's mark. In a rising market,
prices dip at the month's opening and advance at month-end; in a falling
market, they rise at month-opening and fall at month-end---this pattern
continuing for two months or more, up to roughly six months. For methods
of observing this, refer to Upper Volume Chapter~2, Chapter~14, and
Chapter~40.

The range-bound market, on the other hand, occurs in many forms:

\begin{enumerate}
\item A range-bound decline within an overall rise from floor to ceiling;
\item A range-bound rally within an overall fall from ceiling to floor;
\item Range-bound rises and declines within a small-amplitude
  consolidation.
\end{enumerate}

For methods of reading the range-bound market, see Upper Volume
Chapter~10, where the subject is discussed in detail.

% --- Article 61 ---
\article{61}{On the Alignment of Rice, Others, and Self}

\begin{sokyubox}
{\small\itshape\color{sokyubar} \jp{第六十一章\quad 米人我三つ揃候時の事}}

\medskip
When the rice itself is weak, the general sentiment is uniformly
bearish, and one's own assessment is likewise bearish---after such a
time, prices will inevitably rise. When everyone's sentiment is bullish
and the rice appears set to advance, it will certainly fall. As a
general rule, one should regard rice that has risen a hundred bales
(\jt{百俵}{hyappy\=o}) as bound to fall; and rice that is feebly weak
beyond all reason (\textit{tawahi mo naku}) as bound to rise. This is
essential.
\end{sokyubox}

\commentary
The meaning of this chapter is that when the market approaches a ceiling
or floor, both bulls and bears tend to lean uniformly in one direction,
and one is easily swept along by the prevailing sentiment. But one
should not be deceived by the crowd. Using the hundred-bale rise as
one's standard, one should gauge the position of the ceiling or floor,
and take the opposite side of the prevailing sentiment---selling into
strength or buying into weakness. For further reference, see Chapter~42.

% --- Article 62 ---
\article{62}{On Accumulating Buys at the Floor}

\begin{sokyubox}
{\small\itshape\color{sokyubar} \jp{第六十二章\quad 底直段買重ねの事}}

\medskip
When rice has been consolidating at the floor price and begins to rise,
it will reach the ceiling price rapidly within two or three months. At
such a time, using the hundred-bale rise as one's target, it is good to
accumulate additional buy positions (\jt{買重ね}{kai-gasane}).
\end{sokyubox}

\commentary
When rice has been quietly consolidating at the floor for a prolonged
period, building latent momentum, and some catalyst emerges that triggers
an upward breakout, the market will continue rising for two or three
months. At such a time, one should estimate the upward range and,
reckoning on roughly a hundred-bale rise, accumulate buy positions.

A floor-level consolidation typically lasts around thirty days at the
major floor; the usual pattern is for the market to break above the
prior month's monthly mark within forty to fifty days. However, after a
bountiful harvest year or a violent swing, a prolonged consolidation may
sometimes extend to about a hundred days, so one should take note.

% --- Article 63 ---
\article{63}{On How Falling and Rising Markets Differ Greatly}

\begin{sokyubox}
{\small\itshape\color{sokyubar} \jp{第六十三章\quad 下相塲は上相塲と大違の事}}

\medskip
A falling market, from the month the ceiling price is established after
a rise of more than a hundred bales, declines gradually over five or six
months. It differs greatly from a rising market. In terms of the rice
market's oscillation pattern, prices rise at the month-opening and fall
at month-end; each month they may gain roughly one bale at the opening,
then lose roughly two bales by the close. If from the first through the
fourth or fifth of the month, rice drops one or two increments, this
represents a buying opportunity within the range-bound oscillation.
However, since this is ultimately a rally within a larger decline, it is
temporary, and one cannot buy in excess; the position should be closed
out within the month.
\end{sokyubox}

\commentary
This chapter describes the normal behavior and oscillation pattern of a
falling market.

A falling market characteristically shows only modest range-bound
rallies and declines steeply for the most part. The standard pattern is
for prices to rise at the month-opening and fall at month-end. However,
if this pattern is broken and the market shows an unexpected decline at
the very beginning of the month, this signals a range-bound oscillation;
one should buy at the dip. The market will indeed rally by mid-month.
But since this is only a range-bound rally within a larger decline and
not a genuine trend advance, one should keep one's position small.

% --- Article 64 ---
\article{64}{On How Rice Follows the Natural Order, Beyond the Reach of Calculation}

\begin{sokyubox}
{\small\itshape\color{sokyubar} \jp{第六十四章\quad 米は天性自然の理算用に及ばざる事}}

\medskip
When one trades by comparing prices against the Kamigata market
(\jt{上方相塲}{Kamigata soba}) and calculating the volume of arbitrage
rice (\jt{乘米}{nosei-mai}), one goes against the market and invites
error. Rice rises and falls according to the natural order of heaven;
calculation cannot reach it. The ultimate secret is to understand the
sequence of ceiling and floor, and to pay attention to the market every
day. This is essential.
\end{sokyubox}

\commentary
To disregard whether the current price stands at the ceiling or the
floor, and to trade solely on the basis of whether the arbitrage spread
is favorable or unfavorable against other regions, or to maneuver one's
trades purely by such calculations, is a mistake. This means failing to
follow the broad trend and going against the market. The rise and fall of
rice prices is governed by forces that inevitably shift; arithmetic alone
cannot capture it. The ultimate secret lies in attending to the ceiling
and the floor, and to the waxing and waning of market sentiment.

% --- Article 65 ---
\article{65}{The Secret Tradition of Turning One's Sentiment}

\begin{sokyubox}
{\small\itshape\color{sokyubar} \jp{第六十五章\quad 米商の秘傳氣を轉ずる事}}

\medskip
After a major move up or down, when the market enters a consolidation
(\jt{保合}{mochai}), if it looks somehow weak and, no matter how one
considers it or compares with other regions, no bullish signs
appear---and it seems that prices must surely fall another twenty or
thirty bales, and that unless one sells here one will miss the
chance---when the urge to sell presses irresistibly, at that moment,
turn one's sentiment and take the buy side. This is the Tradition.

To take the buy side when one is deeply bearish (\jt{弱氣}{yowaki}) is
extremely frightening; yet if one follows this text, abandons one's own
opinion entirely, and buys, profit is assured. Likewise, when the urge
to buy surges irresistibly, turn one's sentiment and sell. This is the
secret tradition (\jt{秘傳}{hiden}) of the rice trader. One must never
relax this mental discipline. When one feels bullish, one should
recognize that others feel bullish too. The same holds when one feels
bearish.
\end{sokyubox}

\commentary
This chapter describes how to form one's assessment during a
consolidation that follows a major price swing. Its meaning is as
follows: a consolidation market has the effect of making people gloomy;
everyone's sentiment turns bearish, and one's own mood follows suit. It
begins to look as though prices must surely drop another twenty or thirty
increments in the near future, and unless one sells now there will be no
opportunity. The urge to sell becomes intense. But in such a case,
everyone's mind is the same as yours; those who were going to sell have
already sold; the market, along with sentiment, has already declined as
far as it can, exhausting all bearish energy. To sell at such a point is
the height of folly. Rather, one should buy against the prevailing
sentiment.

To buy when one is bearish is a maneuver that is exceedingly difficult
to execute. Yet this is precisely the moment to cast aside one's own
opinion utterly, perceive the signs of a counter-move in sentiment, and
buy. Profit will follow. And when bullish sentiment is surging, one
should apply the same principle and sell against it. This is the secret
tradition of the rice trader. One must never, ever follow the crowd's
sentiment at a turning point.

% --- Article 66 ---
\article{66}{On the Error of Selling While Holding Bullish Conviction}

\begin{sokyubox}
{\small\itshape\color{sokyubar} \jp{第六十六章\quad 買氣をはさみ賣方心得違の事}}

\medskip
To sell while harboring bullish conviction (\jt{買氣}{kaiki}) is a
grave error of judgment. For example, suppose one thinks: ``This rice
will rise, but there is unlikely to be any strength until around
month-end; I will sell a little here, take the interest, and buy back
around the first or second of the next month.'' To sell on such
reasoning is a serious mistake. Invariably, prices begin to rise just
when one expects them to fall. When one has sold in the conviction that
prices must surely decline, and then the market rallies, one finds it
even harder to buy back; one day passes, then two, and one falls behind,
missing the chance to buy---and ends up on the sell side in the very
month one expected to buy. Even if month-end brings a slight dip as
expected, one thinks: ``Since the rice still looks a little weak, surely
the opportunity to take profits by next month will be no different''---and
the weaker rice appears, the more one delays buying. When one has judged
that next month will bring a rise, one should simply wait for the buying
opportunity.
\end{sokyubox}

\commentary
This chapter is a warning against selling for short-term profit when
one's fundamental conviction is bullish.
